\documentclass[11pt,a4paper]{article}

\usepackage[utf8]{inputenc}
\usepackage[T1]{fontenc}
\usepackage[french,english]{babel}
\usepackage{amsmath,amssymb,amsthm}
\usepackage{geometry}
\geometry{margin=2.5cm}
\usepackage{hyperref}

\title{A Modal Dialogue between the Projective Dynamic Logo (PDL) and the Theory of Objectivity (TO)}
\author{Cédric Laubscher\\
\small Independent researcher in theoretical and mathematical physics, Switzerland}
\date{February 2026}

\begin{document}
\maketitle

\begin{abstract}
This article develops a systematic comparative analysis between Projective Dynamic Logo (PDL)—a relational framework that endeavours to reconstruct physical reality from minimal axioms imposed on signed graphs—and the Theory of Objectivity (TO), a modal–axiomatic ontology grounded in the Seven Absolute Truths. We identify structural compatibilities, points of tension, and decision criteria governing the mutual interpretability of the two frameworks. In particular, we refine the ontological treatment of persistence, observation, and boundaries, and investigate how constants and effective interactions can be interpreted as coherence ratios within PDL when constrained by the modal discipline of TO. The objective is not to assert the completeness of either programme, but rather to delineate where logical necessities may be sharpened and where empirically testable constraints can be formulated.
\end{abstract}

\section{Introduction and Aim of the Dialogue}

The Projective Dynamic Logo (PDL) is a foundational research framework proposed by an independent researcher with an engineering background who is currently teaching in a technical institute, and is disseminated openly within the scientific community. It is conceived as a minimalist formalism aimed at elucidating how physical reality might emerge from purely relational structures. Beginning from an initially pre-metric logical field, PDL aims to determine the minimal set of relational constraints under which a persistent, empirically recognisable phenomenological domain can arise. Within this approach, what is regarded as fundamental is neither matter nor spacetime, but rather the capacity to maintain a reproducible distinction through the organisation of relations on finite signed graphs \cite{Laubscher_PDL_long,Laubscher_LDP_fr}.

PDL develops its framework in an upstream, foundational direction. It begins by analysing the logical preconditions under which any entity can exhibit temporally persistent existence in the absence of a prior metric structure or ontological primitives, and then formalises these preconditions as four axioms: minimal binary pulsation, triangular coherence, minimal completeness, and logical optimisation. Subject to these constraints, a first minimal stationary closure arises, namely the complete signed graph on four vertices and six edges, denoted \((4,6)\). This structure is interpreted as a logically stationary regime and is proposed as a prototypical representation of an electron at rest. Building on this, PDL constructs hierarchical composite architectures, in particular for the proton, and advances a reinterpretation of certain physical constants—such as the fine-structure constant, Boltzmann’s constant, and an effective gravitational coupling—as ratios of internal coherence associated with these discrete structures \cite{Laubscher_PDL_teaser,Laubscher_PDL_IMRad}. In this way, PDL does not merely seek to reformulate existing physics in an alternative idiom, but to demonstrate how familiar particles, constants, and interactions can be derived as consequences of a minimal relational logic.  

The Theory of Objectivity (TO), developed by Cabannas, Silva, and collaborators, adopts a related yet distinct approach. It articulates a modal–axiomatic ontology grounded in the Seven Absolute Truths, imposing the requirement that any admissible universe must satisfy specific logical necessities in order to exist in a temporally persistent manner. TO places particular emphasis on the ontological significance of boundaries, the necessity of multi-observer stabilisation for full-fledged existence, and the postulation of a transcendent substrate beyond the existential quantum, and constructs a deductive cosmological framework under these constraints \cite{Cabannas_Silva_TO_book,Cabannas_Silva_TO_neutrinos}. Although PDL and TO differ in their primitive notions and technical formalisms, they share a common objective: to derive the structure of reality from conditions of logical and modal possibility, rather than to posit physical entities and laws as primitive, irreducible givens.

The present article is conceived as a structured dialogue between these two programmes. Its objectives are threefold. First, it aims to map compatibilities, tensions, and decision criteria between PDL’s relational axioms and TO’s Seven Absolute Truths, thereby elucidating the ontological commitments of PDL with respect to existence, persistence, observation, boundaries, and the transcendent. Second, it seeks to reinforce the internal methodological discipline of PDL by explicitly distinguishing results that are derivable from its axioms from those that remain conjectural or parametrically calibrated, and by formulating concrete programmes for mathematical refinement. These include, among others, the problem of the uniqueness of the minimal $(4,6)$ closure, the development of a combinatorial definition of coherence leakage, and a more rigorous treatment of active surfaces and constants. Third, it aims to identify operational interfaces and falsifiable signatures through which the two frameworks can be jointly confronted with empirical and structural constraints \cite{TO_community_about}.  

From the author’s perspective, PDL yields a transparent emergent representation of reality as a single logical totality structured by pulsating relations, within which conventional “domains” such as microphysics, gravitation, and thermodynamics arise as distinct manifestations of an underlying coherence structure. Nonetheless, a substantial portion of this picture remains to be articulated in a form that satisfies prevailing standards of logical, mathematical, and physical rigor. The dialogue with TO is therefore not intended as the presentation of a completed theory, but rather as an opportunity to sharpen the formulation of the emergent logic that grounds PDL, to recast it in terms of explicit axioms and derivations, and to subject it to an externally motivated modal discipline. Philosophical considerations will be invoked only insofar as they contribute to clarifying what it means, within such a minimalist emergent framework, for an entity to exist, to persist, and to generate exportable relational content beyond a single closure.

\section{PDL Axioms and TO's Seven Absolute Truths}

\subsection{The Four Relational Axioms of PDL}

Within the PDL framework, reality is modeled as a network of minimally closed logical structures represented by finite signed graphs. Vertices encode informational entities, while edges encode binary relations of agreement or disagreement. The four axioms formulated below are not proposed as heuristic guidelines, but rather as minimal conditions for the emergence and maintenance of temporally persistent distinctions within an initially pre-metric logical field \cite{Laubscher_PDL_long,Laubscher_LDP_fr}. Where relevant, we state explicitly whether a given axiom is construed as a candidate for modal necessity or as a methodological postulate that remains open to subsequent revision.

\paragraph{Axiom 1 (Minimal binary pulsation).}
Existence is postulated to manifest as an elementary alternation between two mutually exclusive states, denoted abstractly as agreement and disagreement, which is reproducible and definable independently of any antecedent metric structure. Formally, there exists a discrete dynamical system
\[
F : E \times \{\pm 1\} \to E \times \{\pm 1\}
\]
acting on the edge-sign assignments of a finite graph, such that there exists at least one sign configuration \(\sigma\) for which
\[
F(\sigma) = \sigma', \qquad F(\sigma') = \sigma,
\]
with \(\sigma' \neq \sigma\). The minimal invariant associated with existence is thus identified with a logical 2-cycle in the space of sign configurations, independently of any embedding into a continuous space–time framework \cite{Laubscher_PDL_IMRad}. In the current formulation of the framework, Axiom~1 is advanced as a necessity claim: in the absence of such a repeatable alternation, no persistent distinction can be sustained, and any putative ``existence'' becomes logically indistinguishable from nonexistence.

\paragraph{Axiom 2 (Triangular coherence).}
Before formulating Axiom~2, it is useful to justify the choice of triangles as the elementary motifs of closure. In a graph-theoretic framework, a cycle of length three is the shortest possible cycle and constitutes the first configuration in which a distinction can be maintained without reference to any vertex or edge external to the cycle. A single relation, or even a pair of relations, cannot support a reproducible reference without implicitly presupposing an external background; by contrast, a triangle already furnishes a mutually closed pattern, in which each entity is related to the other two and the coherence of the configuration can be assessed locally. In the theory of signed graphs and structural balance, triangles likewise occupy a privileged position: a complete signed graph is balanced if and only if all of its triangles are balanced, i.e.\ the product of the signs along every triad is positive \cite{Harary_balance,Cartwright_Harary}. Within PDL, Axiom~2 can thus be interpreted as adopting, in a minimal manner, this triadic notion of closure as the fundamental unit of logical coherence, rather than positing longer cycles or higher-order motifs as primitive.

A single informational entity, even when equipped with a binary pulsation, is insufficient to sustain a stable and determinate distinction in isolation. Relation is therefore taken to be constitutive. Logical closure is modelled by cycles of length three. Given a finite graph \(G = (V,E)\) with edge signs \(s_{ij} \in \{\pm 1\}\), a triangle \((i,j,k)\) with edges \((i,j), (j,k), (i,k) \in E\) is defined to be \emph{coherent} if
\[
s_{ij} s_{jk} s_{ik} = +1,
\]
and \emph{violated} if at least one of the three edges is absent or if the above product equals \(-1\). Let \(\nu\) denote the fraction of violated triangles,
\[
\nu = \frac{\text{number of violated triangles in } G}{\binom{|V|}{3}},
\]
with the convention \(\nu = 0\) if \(|V| < 3\). Axiom~2 asserts that any structure capable of sustaining a persistent distinction must minimise \(\nu\): violations of triangular coherence diminish stability and undermine closure. Within PDL, this axiom is treated as a candidate for a modal necessity: in the absence of some minimal closure constraint of this form, relational organisation cannot support reproducible reference.

\paragraph{Axiom 3 (Minimal completeness and elementary stationary structures).}
Closure must not depend on any implicit external background. An \emph{elementary stationary structure} is defined as a finite graph \(G = (V,E)\) with signed edges \(s_{ij}\) satisfying the following conditions:
\begin{itemize}
  \item \(G\) is connected and complete, i.e., every pair of distinct vertices is linked by an edge;
  \item there exists an assignment of edge signs such that all triangular subgraphs are coherent (\(\nu = 0\));
  \item the union of all coherent triangles induces a connected subgraph that contains at least two triangles sharing a common edge (hence there is no isolated coherent cycle);
  \item all vertices are structurally equivalent, implying the absence of any imposed internal hierarchy;
  \item the discrete dynamics on the edge signs admits a global binary pulsation \(s \mapsto -s\) that leaves triangular coherence invariant.
\end{itemize}
An elementary stationary structure is said to be minimal if no proper subgraph satisfies all of the stipulated conditions. Within this framework, the first admissible minimal structure occurs for \(|V| = 4\), giving rise to the complete signed graph on four vertices and six edges, denoted by \((4,6)\) \cite{Laubscher_LDP_fr}. Axiom~3 encodes a strong completeness condition: all references required to specify and preserve the distinction must be internal to the structure itself. In the present formulation, this is treated as a mixed assertion: completeness is taken to be a necessary condition for fundamentality, whereas certain aspects of the definition (for example, the exact role of vertex-equivalence) are left open to further refinement.

\paragraph{Axiom 4 (Logical optimisation).}
Within the class of admissible configurations, only those that optimise a prescribed trade-off between coherence, connectivity, and minimality are expected to be stable. To formalise this selection principle, we associate to each configuration \(\sigma\) the following quantities:
\begin{itemize}
  \item \(\nu(\sigma)\), the proportion of violated triangles, as defined in Axiom~2;
  \item \(\rho(\sigma) = r(\sigma)/r_{\max}(|V|)\), the normalised relational density, where \(r(\sigma)\) denotes the number of active edges and \(r_{\max}(n) = n(n-1)/2\) is the number of edges of the complete graph on \(n\) vertices;
  \item a minimality indicator \(m(\sigma) \in \{0,1\}\), which is equal to 1 if \(\sigma\) realises an elementary closed structure without requiring any strictly richer substructure, and 0 otherwise.
\end{itemize}
We introduce an optimisation functional
\[
\Omega(\sigma) = w_{\text{coh}}\,f_{\text{coh}}(\nu(\sigma)) + w_{\text{conn}}\,f_{\text{conn}}(\rho(\sigma)) + w_{\text{min}}\, m(\sigma),
\]
where \(w_{\text{coh}}, w_{\text{conn}}, w_{\text{min}} > 0\) are fixed weights, and \(f_{\text{coh}}, f_{\text{conn}}\) are monotone functions selected so that \(\Omega\) is strictly decreasing in \(\nu\), strictly increasing in \(\rho\) (up to closure), and assumes higher values when minimality is realised \cite{Laubscher_PDL_long}. At the elementary level, the \((4,6)\) structure arises as the first local maximum of \(\Omega\). Axiom~4 is, at this point, explicitly methodological: it codifies the hypothesis that persistence favours local optima of coherence subject to internal closure constraints. A central aspect of the ensuing dialogue with TO will be to determine whether such an optimisation principle can be elevated to the status of a modal necessity claim, or whether it should instead be regarded as a structural conjecture, to be empirically tested and, if necessary, revised.

\subsection{A Brief Recall of TO's Seven Absolute Truths}

In what follows, we do not aim to reconstruct the full formal and conceptual apparatus of the Theory of Objectivity (TO), whose detailed presentation can be found in \cite{Cabannas_Silva_TO_book,Cabannas_Silva_TO_neutrinos,Theory_Objectivity_Web}. Our aim is instead to provide a concise and systematic restatement of the seven modal axioms that constitute its logical nucleus, thereby making the subsequent correspondence with Propositional Dynamic Logic (PDL) explicit and formally unambiguous.

To render the ensuing analysis precise, we recall here the Seven Absolute Truths that ground the modal framework of TO, adopting the formulations given by Cabannas, Silva, and collaborators \cite{Cabannas_Silva_TO_book,Cabannas_Silva_TO_neutrinos}. These axioms are to be interpreted as necessity claims: they specify constraints that any admissible universe must satisfy in order to exist in a temporally persistent and dynamically coherent manner.

\begin{itemize}
  \item \textbf{A1 (Nothingness as primitive mathematical essence).} Nothingness is posited as a primitive, eternal mathematical essence. Any universe must be definable solely in contrast to this primitive notion of nothingness, and its logical characterisation must not presuppose any background structure that is ontologically or mathematically richer than this minimal concept.

  \item \textbf{A2 (Aura and uniqueness).} Every element is associated with a magnetic field or \emph{aura} that confers ontological uniqueness. This aura acts as a structured field of individuation: elements are not merely indexed by arbitrary labels, but are distinguished via specific, intrinsic configuration patterns that cannot be reduced to purely conventional or externally imposed identifiers.

  \item \textbf{A3 (Infinity as non-element).} Infinity is introduced as a non-element that is required for the logical specification of the universe. It functions as a completion operator, a “beyond-element” necessary for any universe to be conceptually closed as a totality, even if all of its local or internal structures remain finite.

  \item \textbf{A4 (Boundary necessity).} Any two distinct elements necessitate the existence of at least one boundary line between them. Genuine distinctness is therefore inseparable from the presence of boundaries: an element is defined not only by its internal constitution, but also by the frontier that separates it from other elements within the universe.

  \item \textbf{A5 (Multi-observer validation).} An element acquires full ontological existence only if it is observed by at least two distinct observers. Existence in the strong ontological sense presupposes multi-observer stabilisation: neither a single observer nor purely internal self-reference suffices to confer complete ontological status.

  \item \textbf{A6 (Compositional ancestry).} Every element is constituted from antecedent elements. There are no absolutely primitive, ancestrally ungrounded units: any structure that exists within a given universe must be, at least in principle, decomposable into prior elements according to a well-specified compositional scheme or generative hierarchy.

  \item \textbf{A7 (Transcendent substance).} No existential universe can arise without a form of substance that transcends its specific quantum or structurally closed domain. For any given quantum system or closed structural totality, there must exist a type of exported substance—often interpreted in TO as informational content generated in atomic relations—that exceeds the internal closure of individual elements and mediates cross-domain coherence.

\end{itemize}

Collectively, these Seven Absolute Truths underpin TO’s commitment to modal necessity, the foundational role of boundaries and closure, the requirement of multi-observer stabilisation for ascribing robust existence, and the postulation of a transcendent substrate beyond any particular quantum domain. In the subsequent subsection, we will analyse how these axioms interface with the relational axioms of PDL, identifying domains of strong compatibility, loci of tension, and the decision conditions that can be formulated as a basis for further research.

Each of these truths is subjected to detailed analysis and formal elaboration in the TO literature, including dedicated logical frameworks and a theorem of universal genesis expressed in TO’s proprietary graph-logical formalism. The compressed formulations provided here serve solely as common reference points for the present discussion and are not intended as an exhaustive exposition of the conceptual content of TO.

\subsection{Compatibilities, Tensions, and Decision Conditions (A1--A7)}

In this subsection, we explicate, for each of TO’s Seven Absolute Truths, the primary points of compatibility with the axioms of PDL, the principal sources of friction highlighted in recent TO-oriented interpretations of PDL, and the associated decision criteria or research programmes for subsequent investigation.

\paragraph{A1 (Nothingness as primitive essence).}
\emph{Compatibility.} PDL adopts an upstream, pre-theoretic stance in which no antecedent metric structure or ontological “furniture” is posited: existence is modelled as the emergence of a stable distinction from an initially undifferentiated logical field. In this respect, the informal notion of a “pre-metric logical field” in PDL may be interpreted as a formalisation of nothingness understood as a primitive mathematical essence, rather than as a substantive ontological background.  
\emph{Tension.} Within TO, a key concern is whether the relational language of PDL is rigorously derived from this primitive nothingness or whether it tacitly presupposes a richer conceptual substrate (e.g., graphs, semiotic structures, or dynamical assumptions) that is introduced without explicit modal or axiomatic justification.  
\emph{Decision condition.} A concrete research programme is to demonstrate that the graph-theoretic apparatus of PDL can be reconstructed from a minimal set of logical constraints on distinguishability and reproducibility, defined relative to a primitive nothingness, rather than being posited as an unconstrained modelling choice.

\paragraph{A2 (Aura and uniqueness).}
\emph{Compatibility.} Within PDL, elements are individuated through relational invariants: the minimal stationary structure \((4,6)\) is characterised by a specific coherent sign configuration and a global pulsation symmetry, while composite structures such as the proton are specified by integer invariants (24, 28, 930, 10\,087, 11\,017) together with an internal organisation into valence, sea, and active surface components. These relational signatures fulfil a role analogous to TO's notion of aura, conceived as a structured field of individuation.  
\emph{Tension.} PDL must determine whether such signatures are purely representational (i.e., artefacts of a particular descriptive framework) or genuinely ontological (i.e., constitutive of the element itself). In TO, uniqueness is not reducible to conventional labels; it must be grounded in necessary structural features.  
\emph{Decision condition.} A central task for PDL is to discriminate which relational invariants are logically forced by the axioms (and by the adopted optimisation principles) and which arise merely as contingent modelling choices. Only invariants of the former type can be legitimately advanced as candidates for an ``aura'' in the TO sense.

\paragraph{A3 (Infinity as non-element).}
\emph{Compatibility.} PDL underscores that local closures coexist within an encompassing relational network. Even when individual structures are finite, their stability and coherence costs depend on an ambient context formed by other closures. This ``beyond-the-closure'' context may be interpreted as playing a role analogous to a non-element required to complete the logical description of the universe.  
\emph{Tension.} PDL currently lacks an explicit operator or axiom corresponding to TO's notion of infinity as a non-elementary completion. Should PDL prove to be logically complete without invoking any notion of an ``outside'', TO will require a robust justification for this omission.  
\emph{Decision condition.} One possible research programme is to formalise, within PDL, a notion of horizon or completion operator associated with families of closures, and to investigate whether such an operator is derivable from the axioms or must instead be introduced as an independent postulate.

\paragraph{A4 (Boundary necessity).}
\emph{Compatibility.} PDL is explicitly formulated around the concept of closed structures and active boundaries. The minimal \((4,6)\) closure realises a first logically identifiable boundary; the proton architecture is characterised by a finite active surface that mediates couplings to an electron-type \((4,6)\) structure; and gravitation is interpreted as a modulation of coherence cost within an emergent metric constructed from compatibility relations among coexisting closures. This framework is strongly aligned with TO’s insistence that genuine distinctions presuppose the existence of boundaries.  
\emph{Tension.} TO treats boundaries as ontologically primitive, rather than as merely computational or representational interfaces. PDL must therefore specify under which conditions a closure is elevated to the status of an “element” endowed with its own boundary, as opposed to remaining a substructure internal to a higher-level closure.  
\emph{Decision condition.} A rigorous PDL criterion is required for determining when a relational closure qualifies as an element (with an associated boundary), particularly in the transition from the \((4,6)\) structure to electrons, and from proton-level closures to nuclei and more complex bound systems.

\paragraph{A5 (Multi-observer validation).}
\emph{Compatibility.} PDL’s coherence constraints are intrinsically multi-relational: triangular coherence already involves triples of entities, and the stability of a closure depends on its compatibility with coexisting closures in the surrounding network. This can be reinterpreted as a form of multi-observation: existence is constrained by the extent to which a regime is corroborated across several relational contexts.  
\emph{Tension.} In its current formulation, PDL accounts for persistence primarily as a self-sustaining stationary regime, without explicitly addressing whether full existence requires relational stabilisation by multiple external closures, as mandated by TO’s A5 condition.  
\emph{Decision condition.} An essential task is to determine whether PDL endorses a strictly internal criterion of existence (closure plus pulsation), or whether it can be generalised to incorporate a requirement of stabilisation by at least two independent environmental closures, thereby yielding a graph-theoretic analogue of TO’s multi-observer condition.

\paragraph{A6 (Compositional ancestry).}
\emph{Compatibility.} PDL constructs composite entities such as the proton as hierarchically organised configurations assembled from more elementary closures: valence cores, relational sea, and active surface. The principle that “every element is composed of elements” is therefore intrinsically instantiated at the levels of architecture and hierarchical decomposition.  
\emph{Tension.} TO imposes a requirement of logical precedence: every element must possess an ancestry in terms of prior elements, and not merely in terms of subordinate substructures. If the \((4,6)\) block is taken as a primitive closure, PDL must specify whether this block is truly without ancestry, or whether its derivation from underlying logical constraints constitutes a form of implicit ancestry.  
\emph{Decision condition.} The dialogue must therefore ascertain whether PDL can articulate a notion of generative step or rule-based ancestry for closures, such that even the minimal \((4,6)\) structure can be regarded as arising from logically prior operations, in a manner consistent with A6.

\paragraph{A7 (Transcendent substance).}
\emph{Compatibility.} PDL introduces a coherence–leakage parameter \(\varepsilon\) associated with composite nucleonic configurations and employs it to construct an effective gravitational coupling and to outline a corresponding cosmological scenario. Under the TO interpretation, in which the transcendent is understood as informational substance exported beyond an individual quantum, \(\varepsilon\) can be regarded as a candidate mechanism for such export: a nonzero fraction of coherence fails to achieve local closure and is instead redistributed across the surrounding relational network.  
\emph{Tension.} In its present formulation, \(\varepsilon\) is specified only in a semi-calibrated fashion and is not yet clearly differentiated, at the ontological level, from ordinary dissipative processes. TO therefore foregrounds the question of whether the transcendent in PDL constitutes a bona fide substance that exceeds the quantum as exportable relational content, or whether it functions merely as a bookkeeping parameter for losses within an otherwise self-contained structure.  
\emph{Decision condition.} A principal research objective is to endow \(\varepsilon\) with a rigorous combinatorial definition as a minimal defect of closure in hierarchical structures, and to elucidate its ontological status as a carrier of exported coherence. Only under these conditions can PDL be said to implement A7 in a manner that is not purely interpretive.


\section{Persistence, Observation, Closure, and the Transcendent}

\subsection{Existence and Temporal Persistence in PDL}

In PDL, the point of departure is a pre-metric framework in which no antecedent concepts of space, time, energy, or substance are postulated. The guiding question is not “what exists?” but rather “under which minimal conditions can something be said to exist in a persistent manner, i.e., to remain distinguishable from nothingness?” In such a framework, a single, non-recurrent event cannot qualify as existence in a strong ontological or operational sense: in the absence of reproducibility, no trace can be sustained, and the putative entity remains logically indistinguishable from non-existence \cite{Laubscher_LDP_fr}.  

The first requirement is therefore the emergence of a minimally reproducible difference. Axiom~1 formalises this requirement as a binary pulsation: existence is instantiated as an elementary alternation between two mutually exclusive states, producing at least one logical 2-cycle on the space of sign configurations. This alternation does not presuppose any temporal metric; it merely stipulates that a difference can be iterated in a logically coherent manner. Axiom~2 then demands that such alternations be embedded within relational structures that minimise violations of triangular coherence. The conjunction of these two axioms yields the notion of a \emph{logical stationary regime}: a configuration whose internal sign dynamics implements a stable pulsation subject to closure constraints, without appeal to any external frame of reference \cite{Laubscher_PDL_long,Laubscher_PDL_IMRad}.

Axiom 3 refines this framework by imposing a condition of minimal completeness: all relational references required to define and preserve a given distinction must be internal to the structure itself. This implies that an elementary persistent structure must take the form of a closed, connected graph in which coherent triangles constitute a covering network and in which all vertices occupy structurally equivalent positions. Under these constraints, the first minimally stationary closed configuration is realised by the complete signed graph on four vertices and six edges, the \((4,6)\) structure. In this sense, PDL associates robust existence with the maintenance of a closed, oscillatory pattern of relations: to exist is to autonomously sustain a self-consistent regime of internal coherence.

Temporal persistence is thus not introduced as an external time parameter, but rather as the capacity of a relational configuration to support an indefinitely repeatable coherence cycle. In subsequent developments of PDL, proper time is reinterpreted as the discrete enumeration of such internal cycles associated with a given stationary regime, while relativistic effects are modelled as modifications in the coherence budget required to sustain these cycles across distinct relational environments. These ideas will be examined in greater detail in later sections; for the purposes of the present discussion, the central thesis is that PDL conceptualises existence and persistence as emergent properties of closed relational dynamics, rather than as primitive attributes of pre-specified objects embedded in an external temporal background.

\subsection{Observation and Multi-Observer Stabilisation}

Within PDL, the axioms are formulated without explicit reference to observers in the standard epistemic sense. Existence and persistence are initially defined as properties of closed relational dynamics: an elementary regime is said to exist if it sustains an internal pulsation subject to the closure constraints imposed by Axioms~1–3. Observation, understood as empirical access through measurement devices or conscious agents, is introduced only at a subsequent stage, when the question arises of how such regimes can be instantiated and interrogated within a concrete physical setting. In this regard, PDL aligns more closely with structural realist or ontic approaches than with epistemic reconstructions: the logical structures are not derived from observational data, but instead provide the pre-existing framework within which observational practices acquire semantic content \cite{Laubscher_PDL_long}.  

Despite this, the relational constraints intrinsic to PDL already encode a substantive form of “multi-validation.” Triangular coherence relations couple three informational entities simultaneously, and the stability of any given closure is contingent on its compatibility with other closures embedded in the surrounding network. A stationary regime is therefore not specified solely by its internal repetitive dynamics; it is additionally characterised by the extent to which its coherence pattern can coexist with, and be corroborated by, neighbouring structures, for example in the construction of composite architectures (protons, nuclei, and more general bound systems). This relational interdependence may be interpreted as a structural analogue of multi-observer confirmation: a regime is viable only to the extent that its internal closure remains consistent across a multiplicity of external relational contexts. 

From the perspective of TO, Axiom~A5 stipulates that an element “fully exists” only if it is observed by at least two distinct others. PDL, in its present formulation, does not yet impose an analogous requirement explicitly at the axiomatic level. The operative criterion of existence is primarily internal (closure combined with pulsation), with external compatibility entering only at the level of composite configurations and the emergent metric structure. A natural direction for future development, prompted by the TO interpretation, is to formalise this dependence on external confirmations more rigorously—for example, by demanding that a stationary closure be regarded as fully real only if its coherence pattern can be embedded into at least two independent, mutually compatible neighbourhoods within the global relational network. Such a condition would furnish a graph-theoretic counterpart to TO’s multi-observer stabilisation principle, without the need to introduce observers as primitive entities into the axiomatic foundation of PDL.

\subsection{Boundaries, Active Surfaces, and the Transcendent Substrate}

In PDL, the concept of closure is intrinsically linked to that of boundary. An elementary stationary structure, such as the \((4,6)\) block, is characterized as a minimal closed configuration: all relational references necessary to establish and preserve its internal differentiation are confined to a finite set of vertices and signed edges. For composite structures, such as the proton, closure is no longer strictly local. Instead, a hierarchical organization emerges, comprising quasi-stationary valence cores, a dynamically evolving relational sea, and a finite active surface through which the composite entity can interact coherently with external closures, in particular with an electron-type \((4,6)\) structure \cite{Laubscher_PDL_IMRad,Laubscher_LDP_fr}.

The active surface functions as a boundary in a strong and technically precise sense. It is defined as the minimal relational envelope that simultaneously (i) separates the internal hierarchical organisation from the external relational environment and (ii) supports the effective coupling channels to other closures. Within the proton model, the active surface is characterised by an effective number of surface relations \(R_{\text{surf}}\), obtained from the valence content via a simple optimisation argument. This finite boundary enables a reinterpretation of the fine-structure constant as a surface-to-volume coherence fraction. In this regard, PDL is closely aligned with TO’s Axiom A4: real distinction is inseparable from the presence of boundaries, and the physically relevant degrees of freedom of a composite structure are concentrated on finite interfaces, rather than being uniformly distributed throughout an undifferentiated bulk \cite{Laubscher_PDL_long}.  

Beyond local boundaries, PDL introduces a coherence-leakage parameter \(\varepsilon\) associated with hierarchical closures. The internal architecture of a proton cannot simultaneously achieve maximal closure at all sublevels (valence cores, sea, surface) and maintain perfect global isolation: a non-zero minimal fraction of coherence necessarily fails to be strictly confined within the composite and must be transferred into the surrounding relational network. This fraction is encoded in \(\varepsilon\), which is then employed to construct an effective gravitational coupling and to outline a cosmological framework in which constraint events in a pre-physical logical field give rise to, and continually recycle, stable closures. Under a TO-inspired interpretation, \(\varepsilon\) is a natural candidate for realising Axiom A7: it can be understood as mediating the export of relational substance beyond the quantum domain, i.e., as a mechanism by which closed structures generate a transcendent contribution to the global relational field \cite{Laubscher_PDL_long,Laubscher_LDP_fr}.

At present, however, \(\varepsilon\) is only partially formalised. In existing PDL formulations, it is introduced as a small but non-zero leakage fraction, calibrated so as to reproduce the empirical order of magnitude of the gravitational constant, and the exponent governing its hierarchical amplification is heuristically motivated rather than derived from first principles. From the perspective of TO, this gives rise to two issues. First, is \(\varepsilon\) merely a bookkeeping device for residual incoherence (a form of dissipation), or does it possess the ontological status of exported substance, in TO’s sense of the transcendent? Second, can \(\varepsilon\) be defined in purely combinatorial terms—as a minimal defect of closure in hierarchical structures selected by the optimisation functional—rather than being tuned a posteriori to fit gravitational data?  

A concrete research programme, and a central objective of the present dialogue, is therefore to endow \(\varepsilon\) with a mathematically precise definition as an invariant of hierarchical closures (for example, in terms of unavoidable coherence violations under global closure constraints), and to elucidate its ontological status. Only under these conditions can PDL legitimately claim to implement A7 as more than an interpretive analogy, and to deliver a genuinely relational account of a transcendent substrate emerging from the systematic failure of perfect closure at composite levels.

\section{Internal Architecture, Constants, and Testability}

\subsection{The Minimal (4,6) Block and Its Status}

Under Axioms~1–3, together with the optimisation principle specified in Axiom~4, the first admissible elementary stationary structure in PDL is the complete signed graph on four vertices and six edges, denoted \((4,6)\). Configurations with \(|V| < 4\) are excluded either because no triangular cycle can be formed (for one or two vertices), or because a single triangle is insufficient to implement a covering network of coherent cycles that remains free of implicit hierarchy (for three vertices). For \(|V| = 4\), by contrast, it becomes possible to construct a complete graph in which: every unordered pair of vertices is connected by an edge; there exists a sign assignment such that all triangular subgraphs are coherent (\(\nu = 0\)); the union of coherent triangles covers the entire graph; and a global sign inversion \(s \mapsto -s\) generates a binary pulsation that preserves coherence \cite{Laubscher_LDP_fr,Laubscher_PDL_IMRad}.  

Within the optimisation framework, the \((4,6)\) structure constitutes the first local maximum of the functional \(\Omega\) at the fundamental level: it saturates relational density (\(\rho = 1\)), minimises violations (\(\nu = 0\)), and realises minimal completeness with all vertices being structurally equivalent. No configuration with fewer entities or fewer relations simultaneously satisfies reproducibility, stability, and autonomy; moreover, any local increase in complexity breaks this minimal character without yielding a commensurate increase in closure. In this sense, \((4,6)\) is not merely a convenient or conventional choice, but the first configuration that fulfils the combined requirements imposed by the axioms and the optimisation principle. 

Within PDL, this structure is advanced as a logical prototype for the electron in its rest frame. The internal 2-cycle \(s \leftrightarrow -s\) is construed as an abstract counterpart of the Compton-frequency oscillation associated with the electron mass, and the reduced Planck constant \(\hbar\) is reinterpreted as the minimal quantum of action associated with a complete coherence cycle of this internal dynamics. The rest energy of the electron thereby appears as the coherence cost required to maintain the elementary stationary regime determined by the \((4,6)\) closure, rather than as an irreducible input parameter \cite{Laubscher_PDL_long}. This identification is deliberately conservative: PDL does not assert that the \((4,6)\) structure *is* the electron in a strong ontological sense, but rather that an electron at rest constitutes the first empirically accessible realisation of the logical stationary regime singled out by the axioms.  

From the perspective of TO, two central questions arise. First, is the emergence of \((4,6)\) as the minimal stationary closure a genuine necessity result under the stated axioms, or does this selection still implicitly rely on modelling choices (e.g., the exclusive emphasis on triangles as closure motifs)? Second, does the correspondence between \((4,6)\) and the electron rest state meet the requirement that necessary properties of the physical entity be reproduced, or does it function primarily as an analogy grounded in pulsation and stability? In the current formulation of PDL, a fully developed mathematical uniqueness theorem for \((4,6)\) has not yet been established, and the proposed correspondence with the electron should be regarded as a structured hypothesis rather than a proven isomorphism. Addressing these issues constitutes a natural decision criterion in the TO–PDL dialogue: either \((4,6)\) can be demonstrated to be uniquely enforced by the axioms, in which case its status as a primitive closure acquires modal robustness, or the space of admissible minimal closures must be broadened and the correspondence with the electron accordingly reformulated.

\subsection{Hierarchical Proton, Active Surface, and the Fine-Structure Constant}

Whereas the electron is represented in PDL as a single elementary stationary closure of type $(4,6)$, the proton necessitates a more elaborate composite characterisation. Empirical findings from quantum chromodynamics (QCD) demonstrate that the proton cannot be adequately modelled as a bound state of three pointlike quarks alone; instead, the dominant contribution to its mass arises from a strongly non-perturbative gluon field and a fluctuating sea of virtual quark–antiquark pairs. Within the PDL framework, this internal complexity is encoded by a hierarchical organisation that integrates multiple locally stationary domains, which are embedded in, and dynamically coupled to, an extensive relational background \cite{Laubscher_PDL_IMRad,Laubscher_LDP_fr}.  

More specifically, the proton is modelled as a two-tier structure composed of three quasi-stationary valence cores, each represented as a superposition of logical tetrads, immersed in a dense relational sea. The valence cores are parametrised by integer quantities $n_u = 24$ and $n_d = 28$ for the up-type and down-type cores, respectively, yielding a total of $R_{\text{valence}} = 930$ effectively stationary relations. These stationary relations are supplemented by a large ensemble of dynamical relations forming a collective sea, with $R_{\text{sea}} \sim 10^4$ active links, so that the total relational content reads $R_{\text{tot}} = R_{\text{valence}} + R_{\text{sea}} \approx 11\,017$. This decomposition indicates that only a small fraction of the relational budget is strictly stationary, while the predominant contribution supports a dynamical regime. The resulting qualitative picture is in agreement with the QCD-based view of the proton as a hierarchically structured and predominantly dynamical system \cite{Laubscher_PDL_long}.  

The hierarchical closure associated with the proton does not couple to an electron-type $(4,6)$ structure through its entire internal volume, but rather via a finite active surface. This active surface is defined as the minimal relational envelope through which the proton can establish coherent interactions with an external elementary closure. In the combinatorial formulation, the active surface is characterised by an effective number of surface relations $R_{\text{surf}}$, derived from the valence content by means of an optimisation argument compatible with Axiom~4. A natural prescription is to regard approximately one third of the valence relations as contributing to an effective core for each of the three quasi-stationary domains, and to multiply this by a dimensionless geometric factor $\varphi$, leading to
\[
R_{\text{surf}} \approx \varphi \,\frac{r_{\text{valence}}}{3},
\]
which yields $R_{\text{surf}} \approx 5.0 \times 10^2$ for $r_{\text{valence}} = 930$ and $\varphi \approx 1.618$.  

The occurrence of the golden ratio $\varphi$ in this expression is not introduced as a mere aesthetic device. From the author’s standpoint, $\varphi$ encodes a geometric constraint intrinsically linked to closure phenomena: when a relational structure that is most naturally specified on a two-dimensional manifold is “folded’’ into a three-dimensional closed configuration, the optimal redistribution of relations between an internal core and an active surface tends to approach $\varphi$-like proportions. 

In discrete terms, this can be formalised by imposing that the proportions of relations allocated to the internal core and to the active surface realise a self-similar partition of the total relational content, such that the functional role of the core with respect to the whole reflects the role of the surface with respect to the core. Under these conditions, one is naturally led to consider ratios of the form
\[
\frac{R_{\text{core}}}{R_{\text{tot}}} = \frac{R_{\text{surf}}}{R_{\text{core}}}
\]
naturally lead to partitions governed by the golden ratio. In the present formulation of PDL, this heuristic is implemented at the level of the proton’s active surface through the factor $\varphi r_{\text{valence}}/3$, although a fully explicit derivation from a clearly specified combinatorial extremum principle remains to be established.\footnote{A more detailed combinatorial analysis of the role of the golden ratio $\varphi$ at the proton surface is provided in the companion note \emph{Emergence of the Golden Ratio in PDL} (Zenodo, 2026).}


A systematic investigation of such self-similar partitions of the relational content, and of the precise conditions under which they yield golden-ratio-like proportions, is deferred to a separate, more technical study.

Within this framework, the fine-structure constant $\alpha$ can be reinterpreted as a surface-to-volume coherence fraction. Let $\mu = m_p/m_e$ denote the proton–electron mass ratio. PDL then motivates an effective expression of the form
\[
\alpha^{-1} \approx \mu \, R_{\text{surf}}^2,
\]
which numerically reproduces the empirical value $\alpha^{-1} \simeq 137$ with a relative deviation of order $10^{-4}$ when $R_{\text{surf}}$ is specified by the above expression. Within this framework, $\alpha$ is interpreted as quantifying the fraction of the proton’s internal coherence that is effectively projected onto its active surface and thereby remains available for electromagnetic coupling, rather than as an unexplained primitive constant. The same relational architecture that guarantees the hierarchical closure of the proton (via its valence and sea structure) and an optimised allocation of relations between core and surface (expressed through the appearance of $\varphi$ and the factor $1/3$ associated with the three valence regimes) simultaneously determines the effective number of surface relations participating in the coupling with an electron-type $(4,6)$.

From the standpoint of TO, several issues naturally arise. First, to what extent is the value of $R_{\text{surf}}$ uniquely fixed by the axioms and optimisation principle of PDL, as opposed to being implicitly tuned to reproduce the observed value of $\alpha$? Second, can the rôle of $\varphi$ be derived from an explicit variational or extremisation problem in the space of admissible surface configurations—formulated in purely combinatorial terms—rather than being introduced as an ad hoc optimisation factor? Third, how stable is the resulting expression for $\alpha$ under perturbations of the underlying combinatorial data (for instance, alternative valence assignments or slightly modified definitions of the active surface)? A systematic analysis of these questions is crucial if the proposed reinterpretation of $\alpha$ is to be regarded as a genuine prediction of minimal closure and optimisation principles within PDL’s relational architecture, rather than merely a numerically accurate but model-contingent reconstruction.

\subsection{Coherence leakage, hierarchical filtering, and gravity}
\label{subsec:coherence_leakage_gravity}

Within the PDL framework, gravitation is not posited as a fundamental field defined on a pre-existing spacetime manifold. Instead, it is construed as the metric manifestation of a minimal \emph{coherence leakage} from composite hierarchical structures into their relational environment. The internal organization of a proton, characterized by $11\,017$ relations distributed between quasi-stationary valence cores and a predominantly dynamical sea, cannot simultaneously fulfill the requirement of maximal closure at all sublevels (valence blocks, sea, active surface) and the requirement of perfect global isolation. Under the optimization functional $\Omega$, there persists an irreducible closure defect: a small fraction of the available coherence cannot be fully confined within the composite and must instead be accommodated by the surrounding network of coexisting closures.

In the terminology of TO, such an exported fraction of coherence constitutes a natural candidate for the instantiation of Axiom A7 (Transcendent substance), provided that $\varepsilon$ is defined combinatorially as a minimal and ineliminable defect of closure, rather than as an empirically introduced dissipation term.

\paragraph{Local definition of the leakage parameter $\varepsilon$.}

To formulate this notion rigorously, we consider the class $G_p$ of signed graphs encoding admissible hierarchical proton architectures, subject to the combinatorial constraints specified in Section~3: valence contents $n_u = 24$ and $n_d = 28$, with $r_{\text{valence}} = 930$ stationary relations; a dynamical sea comprising $R_{\text{sea}} \simeq 10\,087$ relations; a total relational content $R_{\text{tot}} = 11\,017$; a stationary/dynamical partition $(f_{\text{stat}}, f_{\text{dyn}}) \simeq (9\%, 91\%)$; and the presence of an effective active surface $R_{\text{surf}}$ of order $5 \times 10^2$ that mediates coupling to an electron-type $4,6$ block.

For each $G \in G_p$, we introduce an optimisation functional $\Omega(G)$ that aggregates three contributions: the fraction of violated triangles $\eta(G)$, the normalised relational density $\rho(G)$, and a minimality indicator encoding hierarchical economy. We then permit a partial completion of $G$ by adding relations to its environment, under the requirement that $\Omega$ be maximised while all of the above structural constraints remain satisfied. The \emph{minimal combinatorial closure defect of the proton} is defined as
\[
  \varepsilon \;=\; \inf_{G \in G_p} \delta(G),
\]
where $\delta(G)$ denotes the fraction of relational coherence that cannot be confined within the $11\,017$ internal relations of $G$, even after $\Omega$ has been maximised under the proton-level constraints.

In this framework, $\varepsilon$ is a strictly local combinatorial invariant: it quantifies the minimal deviation from perfect closure that any admissible proton hierarchy must exhibit, rather than being introduced as a phenomenological parameter tuned to reproduce gravitational observables. It measures the fraction of relational circulation that must necessarily be exported to the environment, thereby providing a natural candidate for a relational analogue of a ``transcendent'' or externally encoded contribution.

\paragraph{Hierarchical filtering and the exponent 18.}

The same leakage parameter $\varepsilon$, defined at the proton scale, does not map directly onto a macroscopic gravitational coupling. Between the internal hierarchical structure of the proton and the scale at which an effective Newtonian coupling can be operationally inferred, this leakage must propagate through a finite number of independent coherence filters associated with higher organisational levels. From a conceptual standpoint, these filters can be partitioned into three groups of six:

\begin{itemize}
  \item a first block concerned with the selection of an admissible proton architecture, encompassing: the compatibility between the $4,6$ brick and tetrad configurations; the local completeness of $u$- and $d$-type blocks; the minimisation of the $u,u,d$ charge asymmetry; the enforcement of a $9\%/91\%$ stationary/dynamical partition; the existence of a finite active surface; and the reduction of the global closure defect $\varepsilon$ to a minimal value;
  \item a second block concerned with the formation and stabilisation of nuclei, including: proton--proton and proton--neutron compatibility; the closure properties of light nuclei; the allocation of binding degrees of freedom; the preservation of a well-defined leakage per nucleon; and the stability of a discrete nuclear energy spectrum;
  \item a third block concerned with the aggregation of nuclei into ordinary matter and the emergence of a metric regime, involving: nucleus--nucleus compatibility in condensed phases; the transfer of leakage to a mesoscopic scale; the formation of self-gravitating structures; the recovery of an effective $1/r^2$ dependence in the low-leakage regime; the cosmological consistency of residual leakage; and the universality of the effective gravitational coupling.
\end{itemize}

Taken collectively, these requirements define a set of $18$ independent coherence filters that any local leakage, specified at the proton scale, must satisfy in order to be consistently interpreted as a universal, large-scale curvature in an emergent metric.

At the present stage of development, this decomposition should be understood as a structurally motivated counting scheme rather than as a rigorously established uniqueness theorem. The assertion is that any physically admissible mapping from a local leakage parameter $\varepsilon$ to a universal effective coupling $G_{\mathrm{eff}}$ must implement at least this number of logically distinct constraints; it is not claimed that a more fine-grained or higher-resolution description of the underlying structure is precluded. \footnote{A first technical exploration of this hierarchical filtering scheme and of the integer exponent $18$ is presented in the companion note\emph{Hierarchical Coherence Filtering and the Exponent 18 in PDL} (Zenodo, 2026).}

In schematic terms, one may model the effective gravitational coupling as arising from a hierarchical amplification mechanism of the form
\[
  G_{\text{eff}} \;\sim\; \mathcal{C} \,\varepsilon^{\,18},
\]
where $\mathcal{C}$ is a dimensionless prefactor constructed from the same structural elements that enter the definitions of $\alpha$ and $k_B$—in particular, the $4,6$ relations, the proton mass, and the relativistic conversion between coherence cost and energy. The key point is that the exponent $18$ is not introduced as a tunable integer parameter, but is interpreted as the minimal number of level-specific coherence constraints that a single leakage parameter $\varepsilon$ must satisfy in order to:

\begin{enumerate}
  \item define a minimal closure defect compatible with an admissible proton-scale hierarchy;
  \item propagate consistently through nuclear and condensed-matter organisation without being arbitrarily suppressed or amplified;
  \item and give rise to a universal effective coupling $G_{\text{eff}}$ at macroscopic scales, independent of the detailed microscopic chemical composition of baryonic matter.
\end{enumerate}

From this standpoint, both the relative weakness and the universality of gravitation are understood as direct consequences of (i) the small magnitude of the minimal closure defect $\varepsilon$ at the proton level, and (ii) the finite number of hierarchical filters that any such leakage must traverse before being encoded as an effective metric curvature. A full combinatorial derivation of the 18 filters—formulated as explicit conditions on multi-level signed graphs and on the optimisation functional $\Omega$—is deferred to future work. The present discussion is intended primarily to specify their conceptual function and to clarify the interpretative status of both $\varepsilon$ and the exponent $18$ within the PDL framework.

\subsection{Operational Bridges and Falsifiability}

The PDL framework is not proposed as a self-contained replacement for established physical theories, but rather as a research programme that must be formulated in terms of explicit empirical tests and potential refutations. From a TO standpoint, this entails pinpointing those aspects in which the relational axioms and their concrete implementations yield non-trivial, empirically constrainable predictions. Several such “operational bridges” between the abstract formalism and observational or experimental data can be delineated.

At the current stage of development, it is important to maintain a clear conceptual distinction. On the one hand, results concerning the emergence of the $(4,6)$ block and the basic proton structure follow in a comparatively direct manner from the axioms combined with the optimisation principle. On the other hand, the proposed reinterpretation of physical constants and of the gravitational sector still depends on partially conjectural identifications, as well as on auxiliary structures such as the parameter $\varepsilon$ and the exponent $18$, whose full derivation from the underlying combinatorial formulation remains an open task.

First, the emergence of the \((4,6)\) structure as the minimal elementary stationary closure can be reformulated as a precise mathematical and conceptual test. Under Axioms~1–4, PDL yields the prediction that no configuration with a smaller number of entities or relations can simultaneously satisfy the combined requirements of binary pulsation, triangular coherence, minimal completeness, and absence of internal hierarchy. This furnishes a clear route to falsification: either (i) a closure with fewer constituents, or one that is structurally non-isomorphic to \((4,6)\), can be exhibited that meets these constraints, or (ii) it can be rigorously demonstrated that the \((4,6)\) configuration is uniquely implied by the axioms. In the latter scenario, the interpretation of \((4,6)\) as a logical threshold for persistence would be substantially reinforced; in the former, the axiomatic basis would require revision, or the purported necessity of the axioms would have to be reconsidered.

Second, the proposed combinatorial architecture for the proton specifies particular integer-valued parameters and a definite partition between stationary and dynamical relations (24, 28, 930, 10\,087, 11\,017), together with a finite active surface \(R_{\text{surf}} \approx 5.0 \times 10^2\). If these quantities are genuinely obtained as solutions of the optimisation functional \(\Omega\) under explicit constraints (charge, closure, hierarchy), then small, controlled perturbations of the axioms or of the optimisation criteria should induce correspondingly systematic and predictable modifications in the resulting architecture. Conversely, if these integers have been tuned in an ad hoc manner to reproduce known phenomenology, such fine-tuning will be exposed under systematic variation of the axioms and weights. A concrete research programme, therefore, is to investigate the robustness of the proton architecture under such perturbations and to compare the resulting combinatorial patterns with empirical constraints on nucleon structure.

Third, the reinterpretation of fundamental constants—including the fine-structure constant \(\alpha\), Boltzmann’s constant \(k_B\), and an effective gravitational coupling \(G_{\text{eff}}\)—furnishes quantitative links between PDL and empirical observations. The current formulations reproduce the correct orders of magnitude and, in certain instances, relative values at the \(10^{-4}\)–\(10^{-2}\) level, while relying on a single compact set of structural inputs (namely the \((4,6)\) relations, the proton integers, the mass ratio \(\mu\), the geometric factor \(\varphi\), and the leakage parameter \(\varepsilon\)). Refining these correspondences into sharper, testable predictions entails two key steps: (i) eliminating or substantially reducing phenomenological calibration by deriving quantities such as \(R_{\text{surf}}\) and \(\varepsilon\) from purely combinatorial principles, and (ii) assigning quantitative uncertainty estimates and specifying domains of applicability. Subsequent high-precision measurements, or improved combinatorial analyses that produce results lying outside these error bounds, would then constitute a direct means of falsifying the framework.

Fourth, PDL proposes a relational reformulation of proper time (as discrete cycle counting) and of gravitation (as variation in coherence cost within an emergent metric structure). These concepts are, in principle, amenable to confrontation with relativistic phenomenology and astrophysical observations. For instance, the manner in which coherence resources are redistributed within bound systems may impose constraints on the effective gravitational coupling of composite bodies, potentially yielding testable deviations from—or consistency checks with—general relativity in particular regimes. At present, this research direction remains preliminary; the ongoing dialogue with TO indicates that it should be extended into precise, quantitative criteria specifying when and how PDL’s emergent metric must reproduce, approximate, or systematically depart from standard relativistic behaviour.

At a more conceptual level, the TO $\times$ PDL interaction itself operates as an operational test of coherence. If PDL can be systematically and consistently translated into the modal axioms of TO without generating latent contradictions, and if its central constructions (such as minimal closure, hierarchies, constants, and leakage) can be reformulated within the conceptual and syntactic discipline of TO without loss of structural information, this would indicate that PDL satisfies a robust criterion of logical compatibility. Conversely, if essential components of PDL resist such reformulation, this will precisely delineate the points at which its underlying assumptions diverge from TO’s notion of necessity and where subsequent revisions or clarifications become required. 

In this respect, TO does not merely provide an alternative idiom in which to restate PDL; it also functions as an external modal constraint that can exclude certain modelling choices (for example, in the treatment of $\varepsilon$ or of active surfaces) whenever these cannot be justified as necessary, or at least as structurally well-disciplined, within TO’s framework.

In either case, the overall framework remains explicitly open to further refinement and to confrontation with empirical or application-driven considerations, rather than claiming methodological immunity on the basis of purely formal properties alone.

\section{Conclusion: The Role of a TO \texorpdfstring{$\times$}{×} PDL Dialogue}

The objective of this article has not been to advance a complete foundational theory, but rather to establish a rigorously structured dialogue between two research programmes that both attempt to derive physical reality from constraints of logical and modal possibility. On the PDL side, we have recalled how four relational axioms on signed graphs give rise to a minimal elementary closure of type $(4,6)$, to a hierarchical proton architecture comprising valence cores, a relational sea, and an active surface, as well as to a relational reinterpretation of several constants and effective couplings in terms of coherence ratios. On the TO side, we have summarised the Seven Absolute Truths and employed them as a modal constraint framework to assess which components of PDL can plausibly be ascribed the status of necessity, which should presently be regarded as methodological hypotheses, and where additional conceptual and technical clarification is still required.

Within this framework, several structural clarifications become apparent. First, existence and persistence in PDL attain a rigorously defined status as attributes of closed, temporally oscillatory relational configurations, rather than as primitive predicates of objects embedded in an external temporal background. This reformulation is naturally compatible with TO’s requirement that any admissible universe be grounded in conditions that permit a stable distinction to be maintained relative to an undifferentiated nothingness. Second, the significance of boundaries and active interfaces is refined: PDL’s minimal $(4,6)$ block and its hierarchical proton structure, each characterized by finite relational frontiers, furnish explicit realizations of TO’s boundary requirement, while simultaneously illustrating how effective properties such as electric charge and interaction couplings may be localized on interfaces instead of being homogeneously distributed across bulk degrees of freedom. Third, the introduction of a coherence–leakage parameter $\varepsilon$ and its hierarchical amplification via an exponent $18$ is reinterpreted as a combinatorial construction: rather than being empirically tuned to approximate the order of magnitude of $G$, $\varepsilon$ is redefined as a minimal closure defect at the proton scale, and the exponent $18$ is understood as a structurally justified count of independent filtering stages that any leakage must traverse before it manifests as an emergent universal macroscopic curvature.

From the perspective of TO, this dialogue yields explicit decision criteria and research programmes rather than a binary verdict of acceptance or rejection. At least three principal lines of investigation emerge. First, the status of the structure $(4,6)$ as a minimal stationary closure must be examined rigorously: either a full uniqueness theorem is to be established under Axioms~1–4, or alternative closures must be permitted, necessitating corresponding revisions in the associated physical identifications. Second, the combinatorial definition of $\varepsilon$ and the associated decomposition into 18 hierarchical coherence filters must be extended into explicit invariants on multi-level signed graphs and optimisation landscapes, such that the link between local leakage and the effective gravitational coupling no longer depends on implicit calibration. Third, the proposed numerical relations for $\alpha$, $k_B$, and $G_{\mathrm{eff}}$ must be refined into quantitative predictions with specified uncertainties and clearly delineated domains of validity, thereby enabling falsification through more precise measurements or through more advanced combinatorial analysis.

In all these respects, TO operates not as an alternative physical theory, but as an external modal filter that evaluates whether the constructions proposed within PDL can be warranted as necessary, or at least as structurally constrained, rather than merely as heuristic narratives. Conversely, PDL provides TO with a concrete domain in which concepts such as boundaries, aura, and transcendent substance can be rendered in explicit graph-theoretic form and subjected to empirical constraints. The present exchange is therefore best understood as an invitation to treat foundational frameworks not as closed metaphysical systems, but as evolving relational architectures that must continuously mediate among logical necessity, combinatorial structure, and empirical trace.

\bibliographystyle{plain}
\bibliography{pdl_to_dialogue}
% Le fichier .bib pourra contenir :
% - Tes propres travaux PDL/LDP (Zenodo).
% - Les références TO (Cabannas, Silva, Seven Absolute Truths, analyses critiques).
% - Références physiques standard (Harary, QFT, gravitation, etc.).

\end{document}
